\section{结论}\label{sec:conclusion}
围绕吉隆泥石流灾害链，本文建立了从多源环境资料到结构水热机制、时间预测与传播响应的分析体系，主要结论如下。

（1）因素族聚合使关键因素识别同时具有物理含义与统计对照。基于5526个窗口和112项特征，边际异常、条件贡献和相似过程三路线形成土壤水分、空气湿状态、土壤热状态、空气热状态及平均风场的共识TOP5。土壤热通量构成第五席的竞争因素，排序的解释结合变化方向、背景位置和有效重采样分母。

（2）共同参数下的四机制竞争建立了结构、热、水及耦合的条件稳定性图景。2048个设计情景给出稳定裕度、首过阈类别与临界关系；解析热核与细化有限差分的最大有限启动时刻差为0.347小时。温度响应的深度依赖、储排水与损伤的参数补偿，以及141条原始实验的本构对照，为机制区分提供了量化支持。

（3）三类时间预测引擎对应不同的信息条件。其中变形与生存模型以13个外部失稳组和1个右删失站进行分组评价，固定提前24小时发布72小时预测时，逆速度与状态空间方法的完整标签包含分别为3/13和6/13；24/48/72小时及监测尺度比较进一步显示模型表现与发布设计相关。机制预测同时保留已启动、未来启动、不启动和未知质量，使参数条件、概率区间和外部覆盖具有清晰角色。

（4）34个简化区域计算和31个原生计算形成了传播模型竞争体系。以给定释放下的流深、速度、足迹、源项账及目标时间括号为共同输出，二维、孔压、复合固水及夹带过程分别贡献了新的解释量。7个选定运行与4个机械设计关联，得到49条目标记录；全源面ST的同体积格网对照及有界夹带计算展示了几何、离散尺度和新增物源对传播的影响。

（5）跨问联合分析把环境异常、机械启动和释放后的目标响应连成可追溯链条。区域背景、源区温压/变形和沟谷对象响应构成互补监测层，时间与空间尺度、信息可用性及模型敏感性共同支撑防灾信息配置。
\AIcredit
